% p339_1; source PDF page 339; manually verified editable table.
\begin{tikzpicture}[x=1cm,y=-1cm,inner sep=0pt,outer sep=0pt]
\path[use as bounding box] (0,0) rectangle (16.000,9.760);
\draw[line width=.4pt] (0.000,0.000) rectangle (2.250,0.915);
\node[align=center,text width=1.990cm,font={\linespread{1}\fontsize{10}{13.50}\selectfont\bfseries}] at (1.125,0.458) {罪名区分};
\draw[line width=.4pt] (2.250,0.000) rectangle (9.950,0.915);
\node[align=center,text width=7.440cm,font={\linespread{1}\fontsize{10}{13.50}\selectfont\bfseries}] at (6.100,0.458) {区分标准};
\draw[line width=.4pt] (9.950,0.000) rectangle (16.000,0.915);
\node[align=center,text width=5.790cm,font={\linespread{1}\fontsize{10}{13.50}\selectfont\bfseries}] at (12.975,0.458) {要点};
\draw[line width=.4pt] (0.000,0.915) rectangle (2.250,5.845);
\node[align=center,text width=1.990cm,font={\linespread{1}\fontsize{10}{13.50}\selectfont}] at (1.125,3.380) {伪造货币罪};
\draw[line width=.4pt] (2.250,0.915) rectangle (9.950,1.930);
\node[align=left,text width=7.440cm,font={\linespread{1}\fontsize{10}{13.50}\selectfont}] at (6.100,1.422) {方式1：用假材料制作假币};
\draw[line width=.4pt] (2.250,1.930) rectangle (9.950,5.845);
\node[align=left,text width=7.440cm,font={\linespread{1}\fontsize{10}{13.50}\selectfont}] at (6.100,3.887) {方式2：用真币材料制作假币，并且对真币材料完全翻新，不保留原真币的基础。\\{}例1，将真币的碎片作为材料，拼凑成新的假币，属于伪造货币。\\{}例2，将金属货币熔化后，制作成新的假币，属于伪造货币};
\draw[line width=.4pt] (9.950,0.915) rectangle (16.000,5.845);
\node[align=left,text width=5.790cm,font={\linespread{1}\fontsize{10}{13.50}\selectfont}] at (12.975,3.380) {1.只包括正在流通的货币。伪造停止流通的货币并使用的，以诈骗罪论处。（2011年第59题）\\{}2.包括正在流通的境外货币};
\draw[line width=.4pt] (0.000,5.845) rectangle (2.250,9.760);
\node[align=center,text width=1.990cm,font={\linespread{1}\fontsize{10}{13.50}\selectfont}] at (1.125,7.803) {变造货币罪};
\draw[line width=.4pt] (2.250,5.845) rectangle (9.950,9.760);
\node[align=left,text width=7.440cm,font={\linespread{1}\fontsize{10}{13.50}\selectfont}] at (6.100,7.803) {用真币材料制作假币，并且保留原真币的基础。\\{}例1，将硬币周边的金属刮下来，减少硬币金属含量，属于变造。例2，将100元面额的人民币数字改为50元，属于变造。（2013年第14题）\\{}[注意］同时采用伪造和变造的手段，拼凑假币的行为，以伪造货币来论};
\draw[line width=.4pt] (9.950,5.845) rectangle (16.000,9.760);
\node[align=left,text width=5.790cm,font={\linespread{1}\fontsize{10}{13.50}\selectfont}] at (12.975,7.803) {“出售、购买、运输假币罪”“持有、使用假币罪”中的假币，均指伪造的货币，不包括变造的货币（2016年第54题）};
\end{tikzpicture}
